\section{Imaginary-time formalism}\label{app:D.10}

This section supplies the steps behind the imaginary-time results of \cref{sec:4.7}: the antiperiodicity of the Matsubara Green's function and its Fourier series, the occupation as a frequency sum, and the response of independent electrons written as a sum over a product of two Green's functions and evaluated. The residue formulas for Matsubara sums, \cref{eq:C.5-matsum}, and the uniqueness of the analytic continuation are derived in \cref{app:C.5}. The notation is that of \cref{sec:4.7}: operators evolve in imaginary time with $H_\mu = H - \mu_{\mathrm{ch}}\op{N}_e$, $O(\tau) = \ee^{\tau H_\mu}O\,\ee^{-\tau H_\mu}$, and averages are taken in $\op{\varrho} = \ee^{-\beta H_\mu}/\mathcal{Z}$.

\subsection*{Antiperiodicity and the Fourier series}

Take $0 < \tau < \beta$, so that $-\beta < \tau - \beta < 0$. At the negative argument the time ordering in \cref{eq:4.7-Gtau} puts $c^\dagger_{\vect{k}}$ on the left with a minus sign, and with $c_{\vect{k}}(\tau - \beta) = \ee^{(\tau - \beta)H_\mu}c_{\vect{k}}\,\ee^{-(\tau - \beta)H_\mu}$,
\begin{equation*}
  \mathcal{Z}\,\GF(\vect{k},\tau - \beta) = \Tr\bigl[\ee^{-\beta H_\mu}\,c^\dagger_{\vect{k}}\,\ee^{(\tau - \beta)H_\mu}c_{\vect{k}}\,\ee^{(\beta - \tau)H_\mu}\bigr]
  = \Tr\bigl[\ee^{-\tau H_\mu}c^\dagger_{\vect{k}}\,\ee^{\tau H_\mu}\,\ee^{-\beta H_\mu}c_{\vect{k}}\bigr] ,
\end{equation*}
where the cyclic property of the trace moved the last factor $\ee^{(\beta - \tau)H_\mu}$ to the front, and it combined with $\ee^{-\beta H_\mu}$ into $\ee^{-\tau H_\mu}$. Cycling $\ee^{-\beta H_\mu}c_{\vect{k}}$ to the front in the same way gives $\Tr[\ee^{-\beta H_\mu}c_{\vect{k}}\,\ee^{-\tau H_\mu}c^\dagger_{\vect{k}}\,\ee^{\tau H_\mu}]$, and moving $\ee^{\tau H_\mu}$ to the front, where it commutes with $\ee^{-\beta H_\mu}$, gives $\Tr[\ee^{-\beta H_\mu}c_{\vect{k}}(\tau)\,c^\dagger_{\vect{k}}] = -\mathcal{Z}\,\GF(\vect{k},\tau)$. Hence
\begin{equation}
  \GF(\vect{k},\tau - \beta) = -\GF(\vect{k},\tau), \qquad 0 < \tau < \beta .
  \label{eq:D.10-anti}
\end{equation}
For two bosonic operators the time ordering carries no sign, and the same steps give $R_{ab}(\vect{q},\tau - \beta) = +R_{ab}(\vect{q},\tau)$. The traces converge for $-\beta < \tau < \beta$, where $\ee^{-\beta H_\mu}\ee^{\pm\tau H_\mu}$ still decays at large energies, and this is the whole range on which the functions are defined.

A function on $-\beta < \tau < \beta$ that changes sign under $\tau \to \tau - \beta$ has period $2\beta$, and its Fourier series runs over the frequencies $n\pi/\beta = n\pi T$ with integer $n$. The coefficient of $\ee^{-\ii n\pi T\tau}$ is $(2\beta)^{-1}\int_{-\beta}^{\beta}\dd\tau\,\ee^{\ii n\pi T\tau}\GF(\vect{k},\tau)$. Splitting the integral at $\tau = 0$ and shifting the negative half by $\beta$ with \cref{eq:D.10-anti}, it is $(2\beta)^{-1}[1 - \ee^{-\ii n\pi}]\int_0^\beta\dd\tau\,\ee^{\ii n\pi T\tau}\GF(\vect{k},\tau)$. The bracket is $0$ for even $n$ and $2$ for odd $n = 2j + 1$, which leaves the fermionic frequencies $\omega^{\mathrm{F}}_j$ and the pair \cref{eq:4.7-Gmats}. For a periodic function the bracket is $1 + \ee^{-\ii n\pi}$, and only the bosonic frequencies $\omega^{\mathrm{B}}_j$ survive.

\subsection*{The occupation as a frequency sum}

By \cref{eq:4.7-Gmats} at $\tau = -0^+$, $\langle c^\dagger_{\vect{k}}c_{\vect{k}}\rangle = \GF(\vect{k},-0^+) = T\sum_j\ee^{\ii\omega^{\mathrm{F}}_j0^+}\GF(\vect{k},\ii\omega^{\mathrm{F}}_j)$, which is \cref{eq:4.7-occmats}. The factor $\ee^{\ii\omega^{\mathrm{F}}_j0^+}$ cannot be dropped. Since $\GF \to 1/(\ii\omega^{\mathrm{F}}_j)$ at large $j$, the sum without it converges only conditionally, and its symmetric value is $\tfrac12(\GF(\vect{k},0^+) + \GF(\vect{k},0^-))$, the average of the two sides of the jump at $\tau = 0$, which is $\langle c^\dagger c\rangle - \tfrac12$.

To evaluate the sum, let $g(z) = \ee^{z0^+}\GF(\vect{k},z)$ with $\GF(\vect{k},z) = \int\dd\omega'\,\SpecA(\vect{k},\omega')/(z - \omega')$. Near each pole $\ii\omega^{\mathrm{F}}_j$ of $n_{\mathrm{F}}(z) = 1/(\ee^{\beta z} + 1)$ the residue is $-T$ (\cref{app:C.5}), so
\begin{equation*}
  T\sum_j g(\ii\omega^{\mathrm{F}}_j) = -\oint_{C}\frac{\dd z}{2\pi\ii}\,n_{\mathrm{F}}(z)\,g(z) ,
\end{equation*}
with $C$ the small counterclockwise circles of \cref{fig:4-matsubara}a. The function $g$ is analytic off the real axis. In the upper half plane the circles can therefore be merged into the boundary of the half plane, traversed counterclockwise: the line $\omega + \ii0^+$ from $-\infty$ to $+\infty$ and an arc at infinity. In the lower half plane they merge into the line $\omega - \ii0^+$ from $+\infty$ to $-\infty$ and another arc. On the arcs $|g| \sim 1/|z|$, while $n_{\mathrm{F}}(z)\,\ee^{z0^+}$ decays exponentially for $\Real z \to +\infty$ through $n_{\mathrm{F}}$ and for $\Real z \to -\infty$ through the convergence factor, so the arcs contribute nothing. What remains is
\begin{equation}
  T\sum_j\ee^{\ii\omega^{\mathrm{F}}_j0^+}\GF(\vect{k},\ii\omega^{\mathrm{F}}_j) = -\int_{-\infty}^{\infty}\frac{\dd\omega}{2\pi\ii}\,n_{\mathrm{F}}(\omega)\bigl[\GF(\vect{k},\omega + \ii0^+) - \GF(\vect{k},\omega - \ii0^+)\bigr]
  = \int\dd\omega\,n_{\mathrm{F}}(\omega)\,\SpecA(\vect{k},\omega) ,
  \label{eq:D.10-occ}
\end{equation}
where the jump $\GF(\vect{k},\omega + \ii0^+) - \GF(\vect{k},\omega - \ii0^+) = -2\pi\ii\SpecA(\vect{k},\omega)$ follows from the Sokhotski--Plemelj formula \cref{eq:2.4-SP} applied under the integral that defines $\GF(\vect{k},z)$. This is \cref{eq:4.7-occ}, obtained without the Lehmann form. For a free electron, $\SpecA_0 = \delta(\omega - \xi_{\vect{k}})$ and the sum is $n_{\mathrm{F}}(\xi_{\vect{k}})$; equivalently, \cref{eq:C.5-matsum} with the single pole of $g$ at $z = \xi_{\vect{k}}$ gives the residue $n_{\mathrm{F}}(\xi_{\vect{k}})$ directly.

\subsection*{The bubble of independent electrons}

Take free electrons, $H_\mu = \sum_{\vect{k}}\xi_{\vect{k}}c^\dagger_{\vect{k}}c_{\vect{k}}$ with the spin in the label, and $O_a = O_b = n_e(\vect{q}) = \sum_{\vect{k}}c^\dagger_{\vect{k}+\vect{q}}c_{\vect{k}}$ (\cref{eq:2.3-nq}), with $\vect{q} \ne 0$. In imaginary time $c_{\vect{k}}(\tau) = \ee^{-\xi_{\vect{k}}\tau}c_{\vect{k}}$ and $c^\dagger_{\vect{k}}(\tau) = \ee^{+\xi_{\vect{k}}\tau}c^\dagger_{\vect{k}}$, by the steps that gave the real-time form in \cref{sec:2.3}. For $0 < \tau < \beta$, with $n_e(\vect{q})^\dagger = \sum_{\vect{k}}c^\dagger_{\vect{k}}c_{\vect{k}+\vect{q}}$,
\begin{equation*}
  R^0(\vect{q},\tau) = -\frac{1}{\Vol}\sum_{\vect{k},\vect{k}'}\ee^{(\xi_{\vect{k}} - \xi_{\vect{k}+\vect{q}})\tau}\,\bigl\langle c^\dagger_{\vect{k}}c_{\vect{k}+\vect{q}}\,c^\dagger_{\vect{k}'+\vect{q}}c_{\vect{k}'}\bigr\rangle .
\end{equation*}
The Gibbs state of independent electrons is a product over orbitals (\cref{sec:2.5}), and an average of operators is nonzero only if every orbital that loses an electron also gains one. With $\vect{q} \ne 0$ that requires $\vect{k}' = \vect{k}$, and then $\langle c^\dagger_{\vect{k}}c_{\vect{k}+\vect{q}}c^\dagger_{\vect{k}+\vect{q}}c_{\vect{k}}\rangle = n_{\vect{k}}(1 - n_{\vect{k}+\vect{q}})$ with $n_{\vect{k}} = n_{\mathrm{F}}(\xi_{\vect{k}})$, as in \cref{eq:2.5-indep}. The transform is
\begin{equation*}
  \int_0^\beta\dd\tau\,\ee^{(\ii\omega^{\mathrm{B}}_j + \xi_{\vect{k}} - \xi_{\vect{k}+\vect{q}})\tau} = \frac{\ee^{\beta(\xi_{\vect{k}} - \xi_{\vect{k}+\vect{q}})} - 1}{\ii\omega^{\mathrm{B}}_j + \xi_{\vect{k}} - \xi_{\vect{k}+\vect{q}}} .
\end{equation*}
Since $n_{\vect{k}}\ee^{\beta\xi_{\vect{k}}} = 1 - n_{\vect{k}}$ and $(1 - n_{\vect{k}+\vect{q}})\ee^{-\beta\xi_{\vect{k}+\vect{q}}} = n_{\vect{k}+\vect{q}}$, the numerator times $n_{\vect{k}}(1 - n_{\vect{k}+\vect{q}})$ is $(1 - n_{\vect{k}})n_{\vect{k}+\vect{q}} - n_{\vect{k}}(1 - n_{\vect{k}+\vect{q}}) = n_{\vect{k}+\vect{q}} - n_{\vect{k}}$, and
\begin{equation}
  R^0(\vect{q},\ii\omega^{\mathrm{B}}_j) = \frac{1}{\Vol}\sum_{\vect{k}}\frac{n_{\mathrm{F}}(\xi_{\vect{k}}) - n_{\mathrm{F}}(\xi_{\vect{k}+\vect{q}})}{\ii\omega^{\mathrm{B}}_j + \xi_{\vect{k}} - \xi_{\vect{k}+\vect{q}}} ,
  \label{eq:D.10-R0}
\end{equation}
the right-hand side of \cref{eq:4.7-bubble}. Continued to $\omega + \ii0^+$, its imaginary part is $-(\pi/\Vol)\sum_{\vect{k}}[n_{\vect{k}} - n_{\vect{k}+\vect{q}}]\,\delta(\omega - \xi_{\vect{k}+\vect{q}} + \xi_{\vect{k}})$, which by \cref{eq:4.1-ImR} is the difference of the independent-electron transition sums of \cref{eq:2.5-indep}.

The same result comes out of a frequency sum. Consider
\begin{equation*}
  T\sum_{j'}\GF_0(\vect{k},\ii\omega^{\mathrm{F}}_{j'})\,\GF_0(\vect{k} + \vect{q},\ii\omega^{\mathrm{F}}_{j'} + \ii\omega^{\mathrm{B}}_j)
  = T\sum_{j'}\frac{1}{(\ii\omega^{\mathrm{F}}_{j'} - \xi_{\vect{k}})(\ii\omega^{\mathrm{F}}_{j'} + \ii\omega^{\mathrm{B}}_j - \xi_{\vect{k}+\vect{q}})} .
\end{equation*}
The summand falls off as $1/|z|^2$, so no convergence factor is needed, and the fermionic form of \cref{eq:C.5-matsum} applies with $g(z) = [(z - \xi_{\vect{k}})(z + \ii\omega^{\mathrm{B}}_j - \xi_{\vect{k}+\vect{q}})]^{-1}$. Its pole at $z = \xi_{\vect{k}}$ has the residue $n_{\mathrm{F}}(\xi_{\vect{k}})/(\xi_{\vect{k}} + \ii\omega^{\mathrm{B}}_j - \xi_{\vect{k}+\vect{q}})$. Its pole at $z = \xi_{\vect{k}+\vect{q}} - \ii\omega^{\mathrm{B}}_j$ has the residue $n_{\mathrm{F}}(\xi_{\vect{k}+\vect{q}} - \ii\omega^{\mathrm{B}}_j)/(\xi_{\vect{k}+\vect{q}} - \ii\omega^{\mathrm{B}}_j - \xi_{\vect{k}})$, and since $\ee^{-\ii\beta\omega^{\mathrm{B}}_j} = 1$ the Fermi function is $n_{\mathrm{F}}(\xi_{\vect{k}+\vect{q}})$. The sum of the two residues is
\begin{equation}
  T\sum_{j'}\GF_0(\vect{k},\ii\omega^{\mathrm{F}}_{j'})\,\GF_0(\vect{k} + \vect{q},\ii\omega^{\mathrm{F}}_{j'} + \ii\omega^{\mathrm{B}}_j)
  = \frac{n_{\mathrm{F}}(\xi_{\vect{k}}) - n_{\mathrm{F}}(\xi_{\vect{k}+\vect{q}})}{\ii\omega^{\mathrm{B}}_j + \xi_{\vect{k}} - \xi_{\vect{k}+\vect{q}}} ,
  \label{eq:D.10-sum}
\end{equation}
term by term the summand of \cref{eq:D.10-R0}, which proves \cref{eq:4.7-bubble}. Two features recur in every diagram of \cref{sec:4.8}. The bosonic frequency was simplified in the Fermi function before any continuation, $n_{\mathrm{F}}(\xi - \ii\omega^{\mathrm{B}}_j) = n_{\mathrm{F}}(\xi)$; doing it afterwards would give a function of $\omega$ that is not analytic. And the occupation factors appear as differences, $n_{\mathrm{F}}(\xi_{\vect{k}}) - n_{\mathrm{F}}(\xi_{\vect{k}+\vect{q}})$, so that a transition between two occupied or two empty orbitals drops out, which is the Pauli blocking of \cref{sec:2.5}.

\subsection*{The continuation}

\Cref{eq:D.10-R0} is a ratio of finite sums in which $\ii\omega^{\mathrm{B}}_j$ appears only in the denominators. Putting $z$ in its place gives a function that is analytic off the real axis, takes the values \cref{eq:D.10-R0} at every $\ii\omega^{\mathrm{B}}_j$ with $j > 0$, and falls off as $1/|z|$. By the uniqueness argument of \cref{app:C.5}~\cite{baym1961determination} it is the only such function, and $z = \omega + \ii0^+$ gives the retarded response. For data known only numerically at the Matsubara points, the continuation is no longer a substitution, and \cref{app:C.5} describes the difficulty.
